% from ega2-1-fr.tex
\oldpage[II]{73}

Comme le foncteur $r^*$ est exact à droite
(\hyperref[0.4.3.1-fr]{0, 4.3.1}), on déduit de l'homomorphisme
(\hyperref[II.4.1.5.1-fr]{4.1.5.1}), qui est surjectif
(\hyperref[II.4.1.6-fr]{4.1.6}), un homomorphisme surjectif
\[
  r^*(\alpha_1^\sharp):r^*(p^*(\mathcal{E}))\longrightarrow
  r^*(\mathcal{O}_P(1)).
\]
Mais $r^*(p^*(\mathcal{E}))=q^*(\mathcal{E})$, et
$r^*(\mathcal{O}_P(1))$ est localement isomorphe à
$r^*(\mathcal{O}_P)=\mathcal{O}_X$, autrement dit c'est un faisceau
inversible $\mathcal{L}_r$ sur $\mathcal{O}_X$ et on a ainsi défini à partir
de $r$ un $\mathcal{O}_X$-homomorphisme canonique surjectif
\[
\label{II.4.2.1.1-fr}
  \varphi_r:q^*(\mathcal{E})\longrightarrow\mathcal{L}_r.
\tag{4.2.1.1}
\]

Lorsque $Y=\operatorname{Spec}(A)$ est affine, et
$\mathcal{E}=\widetilde{E}$ on explicite encore cet homomorphisme de la façon
suivante~: étant donné $f\in E$, il résulte d'abord de
(\hyperref[II.2.6.3-fr]{2.6.3}) que l'on a
\[
\label{II.4.2.1.2-fr}
  r^{-1}(D_+(f))=X_{\varphi_r^\flat(f)}.
\tag{4.2.1.2}
\]
D'autre part, soit $V$ un ouvert affine de $X$ contenu dans
$r^{-1}(D_+(f))$, et soit $B$ son anneau, qui est une $A$-algèbre~; posons
$S=\mathbf{S}_A(E)$~; la restriction de $r$ à $V$ correspond à un
$A$-homomorphisme $\omega:S_{(f)}\to B$, on a
$q^*(\mathcal{E})|V=\widetilde{E\otimes_A B}$ et
$\mathcal{L}_r|V=\widetilde{L_r}$, où
$L_r=(S(1))_{(f)}\otimes_{S_{(f)}}B_{[\omega]}$
(\hyperref[I.1.6.5-fr]{I, 1.6.5}). La restriction de $\varphi_r$ à
$q^*(\mathcal{E})|V$ correspond au $B$-homomorphisme
$u:E\otimes_A B\to L_r$ qui transforme $x\otimes1$ en
$(x/1)\otimes f=(f/1)\otimes\omega(x/f)$. L'extension canonique de
$\varphi_r$ en un homomorphisme de $\mathcal{O}_X$-Algèbres
\[
  \psi_r:q^*(\mathbf{S}(\mathcal{E}))=
  \mathbf{S}(q^*(\mathcal{E}))\longrightarrow
  \mathbf{S}(\mathcal{L}_r)=\bigoplus_{n\geq0}\mathcal{L}_r^{\otimes n}
\]
est donc telle que la restriction de $\psi_r$ à
$q^*(\mathbf{S}_n(\mathcal{E}))|V$ correspond à l'homomorphisme
$\mathbf{S}_n(E\otimes_A B)=\mathbf{S}_n(E)\otimes_A B\to L_r^{\otimes n}$
qui transforme $s\otimes1$ en
$(f/1)^{\otimes n}\otimes\omega(s/f^n)$.
\end{env}

\begin{env}[4.2.2]
\label{II.4.2.2-fr}
Inversement, supposons donnés un morphisme $q:X\to Y$, un
$\mathcal{O}_X$-Module inversible $\mathcal{L}$ et un
$\mathcal{O}_Y$-Module quasi-cohérent $\mathcal{E}$~; à un homomorphisme
$\varphi:q^*(\mathcal{E})\to\mathcal{L}$ correspond canoniquement un
homomorphisme de $\mathcal{O}_X$-Algèbres quasi-cohérentes
\[
  \psi:\mathbf{S}(q^*(\mathcal{E}))=q^*(\mathbf{S}(\mathcal{E}))
  \longrightarrow\bigoplus_{n\geq0}\mathcal{L}^{\otimes n},
\]
et par suite (\hyperref[II.3.7.1-fr]{3.7.1}) un $Y$-morphisme
$r_{\mathcal{L},\psi}:G(\psi)\to
\operatorname{Proj}(\mathbf{S}(\mathcal{E}))=\mathbf{P}(\mathcal{E})$, que
nous noterons $r_{\mathcal{L},\varphi}$. Si $\varphi$ est surjectif, il en est
de même de $\psi$, donc (\hyperref[II.3.7.4-fr]{3.7.4})
$r_{\mathcal{L},\varphi}$ est partout défini. En outre, avec les notations de
(\hyperref[II.4.2.1-fr]{4.2.1}) et de (4.2.2)~:
\end{env}

\begin{proposition}[4.2.3]
\label{II.4.2.3-fr}
Étant donnés un morphisme $q:X\to Y$ et un $\mathcal{O}_Y$-Module
quasi-cohérent $\mathcal{E}$, les applications
$r\mapsto(\mathcal{L}_r,\varphi_r)$ et
$(\mathcal{L},\varphi)\mapsto r_{\mathcal{L},\varphi}$ mettent en
correspondance biunivoque l'ensemble des $Y$-morphismes
$r:X\to\mathbf{P}(\mathcal{E})$ et l'ensemble des classes d'équivalence des
couples $(\mathcal{L},\varphi)$ formés d'un $\mathcal{O}_X$-Module inversible
$\mathcal{L}$ et d'un homomorphisme surjectif
$\varphi:q^*(\mathcal{E})\to\mathcal{L}$, deux couples
$(\mathcal{L},\varphi)$, $(\mathcal{L}',\varphi')$ étant équivalents s'il
existe un $\mathcal{O}_X$-isomorphisme
$\tau:\mathcal{L}\xrightarrow{\sim}\mathcal{L}'$ tel que
$\varphi'=\tau\circ\varphi$.
\end{proposition}

\begin{proof}
Partons d'abord d'un $Y$-morphisme $r:X\to\mathbf{P}(\mathcal{E})$, formons
$\mathcal{L}_r$ et $\varphi_r$ (\hyperref[II.4.2.1-fr]{4.2.1}), puis
$r'=r_{\mathcal{L}_r,\varphi_r}$~; il résulte aussitôt de
(\hyperref[II.4.2.1-fr]{4.2.1}) et de
(\hyperref[II.3.7.2-fr]{3.7.2}) que les morphismes $r$ et $r'$ sont
identiques (en prenant pour générateur de $\mathcal{L}_r$ l'élément
$(f/1)\otimes1$ pour définir les homomorphismes $v_n$ de
(\hyperref[II.3.7.2-fr]{3.7.2})). Inversement, partons d'un couple
$(\mathcal{L},\varphi)$ et formons

